% Source: https://artofproblemsolving.com/community/c3h3134177_2012_lmt_spring_guts_round_p1p12__lexington_mathematical_tournament
% Stable identity: aops-topic-3134177-q10
Pythagoras gives the other leg length $\sqrt B$, so the condition is $A=\lceil\sqrt B\rceil$. For $A=1$, $C=a_1=2$ and $B=2$, but $\lceil\sqrt2\rceil=2\ne1$. For $A=2$, $C=a_2=5$ and $B=5$, but $\lceil\sqrt5\rceil=3\ne2$. For $A=3$, $C=a_3=3\cdot5-2=13$, so $B=1+3=4$, but $\lceil\sqrt4\rceil=2\ne3$. For $A=4$, $C=a_4=3\cdot13-5=34$ and $B=7$, giving $\lceil\sqrt7\rceil=3\ne4$. For $A=5$, $C=a_5=3\cdot34-13=89$, so $B=17$ and $\lceil\sqrt{17}\rceil=5$. Thus the smallest solution is $\boxed{(5,17,89)}$.
\begin{asy}
usepackage("amsmath");
size(7cm);
pair C=(0,0),A=(sqrt(17),0),D=(0,sqrt(17));draw(A--C--D--cycle);label("$A_0$",A,E);label("$C_0$",C,SW);label("$D_0$",D,N);label("$\sqrt B$",A/2,S);label("$\sqrt B$",D/2,W);label("$\sqrt{2B}$",(A+D)/2,NE);
\end{asy}
